\documentclass[11pt]{article}
\usepackage[T1]{fontenc}
\usepackage[utf8]{inputenc}
\usepackage[french]{babel}
\usepackage{lmodern}
\usepackage{microtype}
\usepackage[a4paper,margin=2.2cm]{geometry}
\usepackage{xcolor}
\usepackage{amsmath,amssymb,amsfonts,amsthm,bm,mathtools}
\usepackage{graphicx}
\usepackage{booktabs,tabularx,array,multirow}
\usepackage{enumitem}
\usepackage{hyperref}
\usepackage{fancyhdr}
\usepackage{titlesec}
\usepackage{tcolorbox}
\usepackage{tikz}
\usepackage{lastpage}
\usepackage{float}
\usepackage{caption}
\hypersetup{colorlinks=true,linkcolor=MidnightBlue,urlcolor=MidnightBlue,citecolor=MidnightBlue,pdfauthor={Yann Sofian Smatti},pdftitle={Segmentation}}
\definecolor{DeepBlue}{HTML}{153B6B}
\definecolor{SoftBlue}{HTML}{EAF2FB}
\definecolor{Accent}{HTML}{6C63FF}
\definecolor{Warm}{HTML}{F2994A}
\definecolor{Forest}{HTML}{1E8E5A}
\definecolor{Charcoal}{HTML}{243447}
\definecolor{MidnightBlue}{HTML}{1F4E79}
\setlength{\parindent}{0pt}
\setlength{\parskip}{0.45em}
\pagestyle{fancy}
\fancyhf{}
\fancyhead[L]{\textcolor{DeepBlue}{\textbf{Segmentation d'images — notes L3 maths}}}
\fancyhead[R]{\textcolor{DeepBlue}{\thepage/\pageref{LastPage}}}
\fancyfoot[C]{\textcolor{Charcoal}{Doctum Consilium — support pédagogique}}
\titleformat{\section}{\Large\bfseries\color{DeepBlue}}{\thesection}{0.6em}{}
\titleformat{\subsection}{\large\bfseries\color{MidnightBlue}}{\thesubsection}{0.5em}{}
\titleformat{\subsubsection}{\normalsize\bfseries\color{Charcoal}}{\thesubsubsection}{0.5em}{}
\titlespacing*{\section}{0pt}{1.2em}{0.6em}
\titlespacing*{\subsection}{0pt}{0.9em}{0.35em}
\newtcolorbox{ideaBox}{colback=SoftBlue,colframe=DeepBlue,boxrule=0.8pt,arc=2mm,left=2mm,right=2mm,top=1.2mm,bottom=1.2mm}
\newtcolorbox{mathBox}{colback=white,colframe=Accent,boxrule=0.8pt,arc=2mm,left=2mm,right=2mm,top=1.2mm,bottom=1.2mm}
\newtcolorbox{warningBox}{colback=white,colframe=Warm,boxrule=0.8pt,arc=2mm,left=2mm,right=2mm,top=1.2mm,bottom=1.2mm}
\newcommand{\R}{\mathbb{R}}
\newcommand{\1}{\mathbf{1}}
\newcommand{\vect}[1]{\bm{#1}}
\newcommand{\coverpage}[3]{%
\begin{titlepage}
\begin{tikzpicture}[remember picture,overlay]
\fill[SoftBlue] (current page.south west) rectangle (current page.north east);
\fill[DeepBlue] (current page.north west) rectangle ([yshift=-4.5cm]current page.north east);
\fill[Accent!85] ([xshift=-3.2cm]current page.south east) rectangle (current page.south east);
\end{tikzpicture}
\vspace*{2cm}
{\color{white}\fontsize{24}{28}\selectfont\bfseries #1\par}
\vspace{0.6cm}
{\color{white}\Large #2\par}
\vfill
\begin{tcolorbox}[colback=white,colframe=white,boxrule=0pt,arc=2mm,width=0.92\textwidth]
\textbf{Objectif.} #3
\end{tcolorbox}
\vspace{0.8cm}
{\large\textcolor{Charcoal}{Version pédagogique, niveau L3 mathématiques / ingénierie ML}}\par
\vspace{0.3cm}
{\textcolor{Charcoal}{Document ecrit par Yann Sofian Smatti}}\par
\end{titlepage}}

\begin{document}
\coverpage{Modèles, pertes et métriques de segmentation}{Du U-Net classique aux variantes modernes}{Ce document présente le pipeline complet d'un projet de segmentation : préparation des données, architectures, pertes, métriques, calibration, seuils, cas multi-tâches et lecture critique des résultats, avec une perspective à la fois mathématique et ingénieur ML.}

\tableofcontents
\newpage

\section{Pipeline général d'un système de segmentation}
Un pipeline de segmentation contient généralement les étapes suivantes :
\begin{enumerate}[label=\arabic*)]
  \item acquisition et annotation des données ;
  \item prétraitement et normalisation ;
  \item choix du modèle ;
  \item choix de la perte ;
  \item entraînement et validation ;
  \item post-traitement ;
  \item évaluation statistique et visuelle.
\end{enumerate}

\begin{ideaBox}
\textbf{Vision d'ensemble.} La segmentation n'est pas seulement un réseau. C'est un problème inverse complet : on transforme une image en une région géométrique crédible et mesurable.
\end{ideaBox}

\section{Prétraitement : pourquoi il change tout}
Pour les images médicales, le prétraitement n'est pas cosmétique. Il affecte directement la distribution statistique des données d'entrée.

\subsection{Normalisation}
Une normalisation simple consiste à centrer-réduire :
\[
\widetilde I = \frac{I-\mu}{\sigma},
\]
où $\mu$ et $\sigma$ sont estimés sur l'image, le batch ou le dataset.

Dans les IRM, on emploie aussi des normalisations robustes fondées sur les quantiles, car les intensités n'ont pas toujours une signification absolue homogène entre machines.

\subsection{Redimensionnement et interpolation}
Si l'on passe d'une grille à une autre, on introduit une interpolation. Cela signifie que l'image fournie au modèle n'est plus exactement l'échantillon brut. C'est une approximation numérique qu'il faut assumer.

\section{Architecture U-Net}
Le U-Net classique alterne blocs de convolution, sous-échantillonnage, puis sur-échantillonnage.

\subsection{Forme abstraite}
On peut écrire l'encodeur comme
\[
E_0=X,
\qquad E_{k+1}=\phi_k(E_k),
\]
où $\phi_k$ combine convolutions, normalisations, non-linéarités et pooling.

Le décodeur reconstruit ensuite
\[
D_{k}=\psi_k(D_{k+1}, E_k),
\]
où la deuxième variable représente la skip connection.

La sortie finale est souvent une convolution $1\times 1$ suivie d'une sigmoid ou d'un softmax.

\subsection{Rôle du pooling}
Le pooling réduit la résolution et augmente le champ réceptif effectif. Il sert donc à intégrer de l'information plus globale.

\begin{mathBox}
\textbf{Champ réceptif.} C'est la région de l'image initiale qui influence une activation donnée dans une couche profonde.
\end{mathBox}

\section{Variantes de U-Net}
\subsection{Attention U-Net}
On ajoute des mécanismes d'attention pour pondérer les zones pertinentes. Formellement, on construit un masque d'attention $\alpha\in[0,1]^{H\times W\times C}$ qui module certaines features :
\[
F_{\text{att}} = \alpha \odot F.
\]

\subsection{U-Net++}
L'idée est d'enrichir les chemins de skip connections par des blocs intermédiaires. On réduit ainsi l'écart sémantique entre encodeur et décodeur.

\subsection{TransUNet / modèles hybrides}
Les modèles hybrides combinent convolutions et attention globale. Les convolutions capturent efficacement les motifs locaux ; l'attention améliore la modélisation des dépendances à longue portée.

\section{Pertes usuelles}
\subsection{Cross-entropie}
En multiclasses, la perte au pixel est
\[
\ell_{ij} = -\sum_{c=1}^C y_{ij}^{(c)}\log p_{ij}^{(c)}.
\]

Elle est bien fondée probabilistiquement, mais elle ne traite pas toujours bien le déséquilibre fort des classes.

\subsection{Dice loss}
On a déjà vu la version soft :
\[
\mathcal{L}_{\mathrm{Dice}} = 1-\frac{2\sum p_{ij}y_{ij}+\varepsilon}{\sum p_{ij}+\sum y_{ij}+\varepsilon}.
\]

\subsection{IoU loss}
De même,
\[
\mathcal{L}_{\mathrm{IoU}} = 1-\frac{\sum p_{ij}y_{ij}+\varepsilon}{\sum p_{ij}+\sum y_{ij}-\sum p_{ij}y_{ij}+\varepsilon}.
\]

\subsection{Focal loss}
Quand les classes sont très déséquilibrées, la focal loss réduit le poids des exemples faciles. En binaire :
\[
\mathcal{L}_{\mathrm{focal}}(p_t) = -\alpha (1-p_t)^\gamma \log(p_t).
\]
Ici $p_t$ désigne la probabilité attribuée à la bonne classe. Le facteur $(1-p_t)^\gamma$ renforce l'apprentissage sur les pixels difficiles.

\subsection{Pertes hybrides}
En pratique, on utilise souvent
\[
\mathcal{L}=\lambda\,\mathcal{L}_{\mathrm{CE}}+(1-\lambda)\,\mathcal{L}_{\mathrm{Dice}}.
\]
L'idée est élégante : la cross-entropie apporte une pression probabiliste locale, la Dice apporte une pression géométrique globale.

\section{Métriques et leur lecture critique}
\subsection{Accuracy pixel}
Elle peut être utile, mais elle est souvent trompeuse si le fond domine.

\subsection{Dice et IoU}
Ce sont les métriques de base pour juger le recouvrement.

\subsection{Sensibilité et précision}
Pour la classe positive,
\[
\mathrm{Sensibilit\'e}=\frac{TP}{TP+FN},
\qquad
\mathrm{Pr\'ecision}=\frac{TP}{TP+FP}.
\]
La sensibilité mesure la capacité à ne pas manquer la lésion. La précision mesure la propreté de la région prédite.

\subsection{Hausdorff distance}
Pour des frontières, on peut comparer les ensembles via une distance géométrique. Si $\partial A$ et $\partial B$ sont deux frontières, la distance de Hausdorff quantifie l'écart maximal entre elles. C'est utile quand la qualité des contours est importante.

\section{Seuil, calibration et post-traitement}
La sortie probabiliste doit être transformée en masque. Dans le cas binaire, on choisit un seuil $\tau$ et on définit
\[
\widehat M_{ij}=\1_{\{p_{ij}\ge \tau\}}.
\]

Choisir $\tau=0.5$ est courant, mais pas toujours optimal. On peut sélectionner le seuil sur l'ensemble de validation pour maximiser Dice ou IoU.

\subsection{Morphologie mathématique}
Après seuillage, on peut appliquer des opérations morphologiques : érosion, dilatation, ouverture, fermeture. Intuitivement, elles permettent de nettoyer des artefacts et de régulariser les formes.

\section{Segmentation multitâche : masque + classe}
Dans certains projets, on souhaite que le même réseau prédise à la fois un masque et une classe globale. On partage alors un encodeur et on ajoute deux têtes :
\begin{itemize}
  \item une tête de segmentation ;
  \item une tête de classification.
\end{itemize}

La perte totale s'écrit par exemple
\[
\mathcal{L}_{\mathrm{tot}} = \lambda_{\mathrm{seg}}\mathcal{L}_{\mathrm{seg}} + \lambda_{\mathrm{cls}}\mathcal{L}_{\mathrm{cls}}.
\]

\begin{ideaBox}
\textbf{Intuition.} La classification pousse l'encodeur à apprendre des signatures globales ; la segmentation l'oblige à conserver la structure spatiale. Les deux tâches peuvent donc s'aider mutuellement.
\end{ideaBox}

\section{Lecture critique d'un résultat expérimental}
Supposons que deux modèles donnent :
\begin{center}
\begin{tabularx}{0.95\textwidth}{>{\bfseries}lcccX}
\toprule
Modèle & Dice & IoU & Sensibilité & Lecture rapide \\
\midrule
U-Net A & 0.89 & 0.81 & 0.93 & Très bon recouvrement, peu de lésions manquées. \\
U-Net B & 0.91 & 0.84 & 0.78 & Meilleure compacité globale, mais risque de manquer des régions. \\
\bottomrule
\end{tabularx}
\end{center}

Le meilleur modèle dépend du contexte. En médecine, manquer une lésion peut coûter plus cher qu'avoir quelques faux positifs supplémentaires.

\section{Lien avec MLflow et DVC}
Dans un projet sérieux, les notions mathématiques doivent être reliées à une traçabilité technique.

\subsection{Avec MLflow}
On loggue :
\begin{itemize}
  \item la configuration ;
  \item la perte ;
  \item Dice et IoU par époque ;
  \item les artefacts visuels : overlays, confusion matrices, exemples de masques.
\end{itemize}

\subsection{Avec DVC}
On versionne :
\begin{itemize}
  \item les datasets bruts ;
  \item les manifests image/masque ;
  \item les modèles ;
  \item les métriques.
\end{itemize}

\begin{warningBox}
\textbf{Bonne pratique.} En segmentation, il faut versionner non seulement les images, mais aussi les masques et les scripts qui les appairent. Sinon, la reproductibilité est illusoire.
\end{warningBox}

\section{Notions avancées à explorer ensuite}
\begin{enumerate}[label=\arabic*)]
  \item segmentation 3D pour IRM volumique ;
  \item supervision faible à partir de Grad-CAM ;
  \item semi-supervisé ;
  \item distillation ;
  \item calibration d'incertitude ;
  \item modèles fondation et promptable segmentation.
\end{enumerate}

\section{Fiche de synthèse}
\begin{center}
\begin{tabularx}{0.98\textwidth}{>{\bfseries}lX}
\toprule
Notion & À retenir \\
\midrule
Image & Fonction discrète sur une grille. \\
Masque & Fonction indicatrice ou label par pixel. \\
Convolution & Opérateur local linéaire partageant ses poids. \\
U-Net & Contexte global + détails locaux grâce aux skip connections. \\
Cross-entropie & Vision probabiliste locale. \\
Dice loss & Vision géométrique globale du recouvrement. \\
IoU & Recouvrement plus sévèrement pénalisé. \\
Seuil & Transforme une carte de probabilités en masque binaire. \\
Hausdorff & Mesure l'écart de frontière. \\
Multitâche & Masque et classe globale dans un même réseau. \\
\bottomrule
\end{tabularx}
\end{center}

\section{Exercices guidés}
\subsection*{Exercice 1}
Pourquoi une combinaison CE + Dice est-elle souvent plus stable qu'une Dice loss seule au début de l'entraînement ?

\subsection*{Exercice 2}
Si une lésion est minuscule, quelle métrique surveillerais-tu en priorité parmi accuracy, Dice, sensibilité ? Justifier.

\subsection*{Exercice 3}
Dans un schéma multitâche, que se passe-t-il si $\lambda_{\mathrm{cls}}$ est énorme et $\lambda_{\mathrm{seg}}$ très petit ?

\subsection*{Pistes de réponses}
Une CE pure force chaque pixel à être probabilistiquement cohérent ; la Dice pure regarde surtout le recouvrement global, ce qui peut produire des gradients moins stables au tout début.

\end{document}
